\section{The setup: parabolic data on the family curve}
% =====================================================================
The family curve is the rigid four-punctured sphere $X_f^\circ\cong\mathbb P^1\setminus\mu_4$,
$\mu_4=\{1,i,-1,-i\}$ (the flavor construction above). The flavor bundle is a rank-$\Nfam=3$ parabolic bundle $E_\bullet$ on
$(\mathbb P^1,\mu_4)$. At each puncture the local monodromy is the order-$3$ cusp class with eigenvalues
$\{1,\omega,\omega^2\}$, $\omega=e^{2\pi i/3}$; by the Mehta--Seshadri dictionary the corresponding
\emph{parabolic weights} are the cusp exponents
\begin{equation}
  \boxed{\ \alpha=\Bigl(0,\tfrac13,\tfrac23\Bigr)\ \text{at each of the four punctures}\ }\qquad\mE ,
\end{equation}
i.e.\ the three transport cusps $\{1,\tfrac23,\tfrac13\}=\{2y_+,-2y_-,2(y_++y_-)\}$ written as
exponents $e^{2\pi i\alpha}$ from the $D_5$ dual roots $y_\pm=\{\tfrac12,-\tfrac13\}$.

\begin{figure}[h]\centering
\begin{tikzpicture}[font=\footnotesize,scale=1]
% --- left: P^1 minus mu_4 ---
\begin{scope}[shift={(0,0)}]
  \draw[thick,blue!55!black] (0,0) circle (1.5);
  \node[blue!55!black] at (0,1.85) {\scriptsize $\mathbb P^1\setminus\mu_4$};
  \foreach \ang/\lab in {0/{$1$},90/{$i$},180/{$-1$},270/{$-i$}}{
     \fill[orange!80!black] (\ang:1.5) circle (2.2pt);
     \node[orange!60!black] at (\ang:1.85) {\scriptsize\lab};}
  \node at (0,0) {\scriptsize $A_3=\mathfrak{su}_4$};
  \node[align=center,below=2pt] at (0,-1.9) {\scriptsize 4 cusps $\Rightarrow$ anchor\\[-1pt]\scriptsize $a=(1,1,2)$, $\sum a=4=|\mu_4|$};
\end{scope}
% --- right: C6 hexagon with sector residues ---
\begin{scope}[shift={(6,0)}]
  \draw[thick,gray!50] (0,0) circle (1.5);
  \foreach \k in {0,...,5}{
     \coordinate (h\k) at ({90-60*\k}:1.5);
     \fill[gray!30] (h\k) circle (1.2pt);
     \node[gray!60] at ({90-60*\k}:1.82) {\scriptsize \k};}
  \draw[gray!40] (h0)--(h1)--(h2)--(h3)--(h4)--(h5)--cycle;
  % sector residue sets: u={0,1,3} d={1,2,5} e={2,3,5}
  \node[blue!60!black] at (h0) {$\bullet$}; \node[blue!60!black] at (h1) {$\bullet$}; \node[blue!60!black] at (h3) {$\bullet$};
  \node at (0,0.15) {\scriptsize $R_u{=}\{0,1,3\}$};
  \node at (0,-0.2) {\scriptsize\textcolor{red!70!black}{$R_d{=}\{1,2,5\}$}};
  \node at (0,-0.55) {\scriptsize\textcolor{green!50!black}{$R_e{=}\{2,3,5\}$}};
  \node[align=center,below=2pt] at (0,-1.9) {\scriptsize $C_6$ transport hexagon\\[-1pt]\scriptsize distinct distances $\Rightarrow\{0,1,3\}$ up to $D_6$};
\end{scope}
\draw[-{Stealth[length=2mm]},thick,gray] (1.7,0)--(4.3,0) node[midway,above,align=center]{\scriptsize corner rot.\\[-2pt]\scriptsize $=A_3$ Coxeter};
\end{tikzpicture}
\caption{\small Left: the rigid four-punctured sphere $\mathbb P^1\setminus\mu_4$ ($\mu_4$ as the four
cusps) giving $A_3$; the corner rotation $z\mapsto iz$ is the $A_3$ Coxeter element. Right: the $C_6$
transport hexagon carrying the three sector residue packets; the distinct-distance lemma forces the base
triple $\{0,1,3\}$ up to $D_6$.}
\end{figure}

% =====================================================================
\section{Exact scaffolding (proved)}
% =====================================================================
\begin{proposition}[Parabolic degree zero fixes the underlying degree]
For a unitary (flat) family local system, $E_\bullet$ is parabolic-polystable of parabolic degree zero
(Mehta--Seshadri). With full flags and weights $(0,\tfrac13,\tfrac23)$ at each of the four punctures,
\begin{equation}
  \pardeg(E)=\deg E+\sum_{p\in\mu_4}\bigl(0+\tfrac13+\tfrac23\bigr)=\deg E+4=0
  \quad\Longrightarrow\quad \boxed{\ \deg E=-|\mu_4|=-4\ }\qquad\mE .
\end{equation}
The balanced (stable) underlying splitting is therefore $E\cong\mathcal O(-2)\oplus\mathcal O(-1)\oplus\mathcal O(-1)$. \mE
\end{proposition}

\begin{remark}
$\deg E=-|\mu_4|$ is the same ``four corners'' integer that gives the $\mu_4$ glue and the $41$ chain:
the underlying flavor bundle has degree exactly minus the number of family corners.
\end{remark}

\begin{proposition}[Word-length splitting: windings are degrees, residues live in $\mathbb Z_6$]
Write each diagonal word-length as $L_{f,j}=6\,n_{f,j}+r_{f,j}$. Then the winding $n_{f,j}$ is the
integer degree of the generation's parabolic line sub-bundle, and the residue is
\begin{equation}
  \boxed{\ r_{f,j}\equiv \underbrace{6\,w_{f,j}}_{\in\{0,2,4\}}+\underbrace{3\,s_{f,j}}_{\text{sheet}}\!\!\pmod 6,\qquad
  w_{f,j}\in\{0,\tfrac13,\tfrac23\},\ s_{f,j}\in\{0,1\}\ }\qquad\mE ,
\end{equation}
because the cusp weights contribute $6\alpha\in\{0,2,4\}$ (even) and the $\mathbb Z_2$ sheet parity (the
deck involution $\tau$) shifts by $3$. The hexagon $\mathbb Z_6=\operatorname{lcm}(3_{\text{cusp}},2_{\text{sheet}})$
is exactly $|R^+(A_3)|$.
\end{proposition}

\noindent The TFPT budget is then bookkept as
\begin{equation}
  \sum_{f,j}L_{f,j}=\underbrace{\sum r_{f,j}}_{=22}+6\underbrace{\sum n_{f,j}}_{=\Nfam=3}=22+18=40=|R(D_5)| ,
\end{equation}
so the three windings are exactly the three family line-bundle degree shifts, and the residue sum $22$
is the non-primitive $\mathbb Z_{30}$ load $h-r$ of the main note. \mE

% =====================================================================
\section{The Riemann--Hilbert side (what the monodromy fixes)}
% =====================================================================
Constructing the rigid $\mathbb Z_4$-symmetric local system explicitly (four monodromies
$M_k=R^kCR^{-k}$, $C\sim\operatorname{diag}(1,\omega,\omega^2)$, $R^4=\mathbb 1$), the closure
$M_0M_1M_2M_3=\mathbb 1$ is solvable to machine precision, and the unique rotation $R$ has eigenvalues
$\{i,-1,-i\}$: the $A_3$ Coxeter element. \mE\ This independently reproduces the corner-rotation
$=$ Coxeter identity and confirms the family connection exists.

\par\smallskip\noindent\textbf{Rigidity: honest count.}
The bare character variety of four regular-semisimple $SL_3$ classes on $\mathbb P^1$ has dimension
$\dim\mathcal M=-2\dim SL_3+4\cdot 6=-16+24=8$. The $\mathbb Z_4$-reduction (four punctures $\to$ one
orbifold image, plus the two rotation fixed points) gives a three-point system, still with
$\dim\mathcal M=-16+3\cdot 6=2$ for three regular classes --- \emph{not} cohomologically rigid on its
own (a hypergeometric/pseudo-reflection point would be needed for $\dim=0$). Full rigidity is therefore
the \emph{full} $D_4$ (including the reflection $z\mapsto 1/z$) together with the determinant-trivial
branch, i.e.\ the flavor classification above, which the explicit solution above is consistent with. We do not
claim Katz rigidity for the bare $\mathbb Z_4$ system.
\par\smallskip

\par\smallskip\noindent\textbf{Why the monodromy alone does not give the integers.}
The flat monodromy fixes the \emph{exponents} $\alpha\in\{0,\tfrac13,\tfrac23\}$ but not the integer
\emph{degrees} $n_{f,j}$ nor the per-generation flag positions. The integers are parabolic degrees ---
Hodge/metric data --- so they require the stable parabolic structure, not just the local system. This is
precisely why the residue assignment cannot be read off from the monodromy and is the content of the
project below.
\par\smallskip

% =====================================================================
\section{The parabolic-weight computation (resolved this round)}
% =====================================================================
\noindent\textbf{Status (no drift).} The splitting type is now \emph{closed}: the unique partition
$4{=}1{+}1{+}2$ fixes the anchor and $\operatorname{coker}R=\Z_8$ fixes the residue branch (verified, below
and \veri{v4\_flavor\_matrix.py}), and the solar texture is built (\veri{v9\_neutrino\_texture.py}). The
roadmap below is kept as the \emph{derivation record} of how $L_{f,j}$ is obtained; it is no longer an open
project. The only residual is the parabolic$\leftrightarrow$transport equivalence, conditional on
Mehta--Seshadri unitarity (U).
\par\smallskip\noindent\textbf{Derivation record for $L_{f,j}$ (how the word-lengths are obtained).}
\textbf{(1)} Fix the $D_4$-equivariant stable parabolic structure on
$E=\mathcal O(-2)\oplus\mathcal O(-1)\oplus\mathcal O(-1)$ (the descended three-orbifold-point bundle
with weights: cusp class $\{0,\tfrac13,\tfrac23\}$ at the puncture image; $A_3$-Coxeter exponents at the
two rotation fixed points $0,\infty$).\\
\textbf{(2)} For each sector $f\in\{u,d,e\}$ select the sector cusp
$\{2y_+,-2y_-,2(y_++y_-)\}=\{1,\tfrac23,\tfrac13\}$ (a $D_5$ datum) as the reference weight; this is the
$D_6$ orbit shift $\{0,+2,-2\}$ on the residue triple.\\
\textbf{(3)} Compute, per generation $j$ (an $R$-eigenline, $R$-eigenvalue $\{i,-1,-i\}$), the parabolic
degree of the corresponding parabolic line sub-bundle: the integer part is the winding $n_{f,j}$, the
fractional part times $6$ (plus the sheet) is the residue $r_{f,j}$.\\
\textbf{(4)} Verify the output equals
$\{0,1,3\}$ (up), $\{1,2,5\}$ (down), $\{2,3,5\}$ (lepton), and the per-sector anchors $(1,1,2)$.
\par\smallskip

\noindent\textbf{What is already pinned down for step (3).} The residue must be of the form
$6w+3s\bmod 6$ with $w\in\{0,\tfrac13,\tfrac23\}$; the candidate compact map of the main note,
$a_f=\operatorname{round}(2\,\mathrm{cusp}_f)=(1,1,2)$, is the predicted first-generation anchor and must
emerge as the minimal-weight (dominant) flag step of each sector.

\begin{proposition}[Single-fibre reduction]
Because each generation is an $R$-eigenline ($Rv_j=\zeta_j v_j$) and the puncture flags are $R$-rotated
($M_k=R^kCR^{-k}$, so $F_\bullet(M_k)=R^kF_\bullet(C)$), the weight of $v_j$ is the same at every
puncture, $\alpha(v_j,M_k)=\alpha(v_j,C)$ for all $k$. Hence the parabolic degree collapses to
\begin{equation}
  \boxed{\ \mu_j=d_j+\textstyle\sum_{k=1}^4\alpha(v_j,M_k)=d_j+4\,\alpha_j,\qquad \alpha_j=\alpha(v_j,C)\ }\qquad\mE .
\end{equation}
\end{proposition}

\par\smallskip\noindent\textbf{What the explicit computation shows: the local model is insufficient.}
Evaluating $\alpha_j$ from the Riemann--Hilbert solution (the $R$-eigenvectors against the $C$-flag,
$F_1{=}\langle e_1\rangle$ weight $0$, $F_2{=}\langle e_1,e_2\rangle$ weight $\tfrac13$, $F_3$ weight
$\tfrac23$) gives \emph{degenerate} weights --- generically all $\tfrac23$, at best $\{\tfrac23,0,\tfrac23\}$
--- not the non-degenerate permutation $\{0,\tfrac13,\tfrac23\}$. This is decisive: the residues are
\textbf{not} a single-fibre (monodromy) quantity. The generations are \emph{global} parabolic line
sub-bundles whose weight at each puncture comes from the holomorphic degeneration of the saturated
sub-sheaf, not from the fibre alignment of an eigenvector. So the computation genuinely requires the
global stable parabolic structure (the holomorphic Hecke modification matching the weights, with
parabolic stability fixing the saturation), confirming the $\mC$ classification: it is an
algebraic-geometry computation on $\mathbb P^1$, not a linear-algebra one on a single fibre.
\par\smallskip

\noindent The open piece is therefore the explicit \emph{global} parabolic-degree evaluation at the
$D_4$-fixed point --- the holomorphic Hecke modifications of $\mathcal O(-2)\oplus\mathcal O(-1)^2$
matching the prescribed weights, saturated for parabolic stability --- not a search and not a free
choice, but a genuine moduli computation. The single-fibre reduction $\mu_j=d_j+4\alpha_j$ shows the
\emph{form} of the answer; the holomorphic structure supplies the integers $d_j$ and the global weights
$\alpha_j$.

% =====================================================================
\section{Global computation: the anchor multiset is derived}
% =====================================================================
The global (Fuchsian/Hecke) computation does resolve one integer datum exactly. Work in the rotation
eigenbasis $R=\operatorname{diag}(i,-1,-i)$ and take the $\mathbb Z_4$-symmetric Fuchsian residues
$A_k=R^kA_0R^{-k}$ with $\operatorname{spec}(A_0)=\{0,\tfrac13,\tfrac23\}$ (the cusp exponents). Then the
symmetrisation collapses exactly,
\begin{equation}
  \boxed{\ \sum_{k=0}^{3}R^kA_0R^{-k}=4\operatorname{diag}(A_0)\ }\qquad\mE ,
\end{equation}
because $\sum_{k=0}^3(r_a/r_b)^k=4\,\delta_{ab}$ for $r=(i,-1,-i)$ (off-diagonal terms are summed fourth
roots of unity). The exponents at infinity --- equivalently the splitting type of the Deligne extension,
$A_\infty=-\sum_k A_k$ --- are therefore $4\operatorname{diag}(A_0)$, integers (regularity at $\infty$)
summing to $4\operatorname{tr}(A_0)=4$.

\begin{theorem}[The anchor multiset from stability]
With $\operatorname{spec}(A_0)=\{0,\tfrac13,\tfrac23\}$ the diagonal entries satisfy
$4(A_0)_{aa}=\tfrac43(q_{a2}+2q_{a3})\in[0,\tfrac83)$, so each exponent at infinity is in $\{0,1,2\}$
(the value $3$ is \emph{unreachable}). By Schur--Horn the only integer multisets of exponents
realisable with this spectrum are
\[
  \{0,2,2\}\quad\text{and}\quad\{1,1,2\}.
\]
Parabolic stability (minimal-variance / balanced splitting; $\{0,2,2\}$ has a spurious weight-$0$ flat
direction) selects
\begin{equation}
  \boxed{\ \text{exponents at }\infty=\{1,1,2\}=(a_u,a_d,a_e)=-\deg\bigl(\mathcal O(-2)\oplus\mathcal O(-1)^2\bigr)\ }\qquad\mE .
\end{equation}
An explicit realising matrix is
\[
  A_0=\begin{pmatrix} \tfrac14 & -\tfrac1{12} & \tfrac{\sqrt6}{12}\\[2pt]
  -\tfrac1{12} & \tfrac14 & \tfrac{\sqrt6}{12}\\[2pt]
  \tfrac{\sqrt6}{12} & \tfrac{\sqrt6}{12} & \tfrac12 \end{pmatrix},\qquad
  \operatorname{diag}(A_0)=\bigl(\tfrac14,\tfrac14,\tfrac12\bigr),\quad \operatorname{spec}(A_0)=\{0,\tfrac13,\tfrac23\}
\]
(verified), so the Schur--Horn point realising the anchor $(1,1,2)$ is explicit, not merely abstract.
\end{theorem}

\noindent So the three per-sector \emph{anchors} $(1,1,2)$ are \emph{not} free: they are the balanced
integer splitting type of the $\mathbb Z_4$-symmetric parabolic flavor connection, fixed by Schur--Horn
plus stability, and they coincide with the underlying bundle degrees $\mathcal O(-2)\oplus\mathcal O(-1)^2$.
The ``balanced'' selection is now \emph{typed}, not heuristic: $\{0,2,2\}$ contains an $\mathcal O(0)$
summand with $h^0=1$, while stability of the Deligne extension (irreducible unitary monodromy,
\vref{v117}) forces $h^0(E)=0$ --- positivity \emph{is} the selection, and the Schur--Horn and
Deligne/Birkhoff--Grothendieck routes converge on $\{1,1,2\}$
(\veri{v499\_p2\_weights\_deligne\_bg.py}, \texttt{P2.TYPING.01}). The \vref{v499} residual R1 (the
module identification) is now identified with the QGEO modulus rest: both reduce to the \emph{same}
order-$4$ carrier input --- the clock $=$ Coxeter of $W(A_3)$ $=$ the cusp class
(\veri{v503\_qgeo\_emergence\_light.py}, \texttt{QGEO.EMERGE.LIGHT.01}; one common residual, markers
unchanged). That common residual is now \emph{one alignment bit} --- the collar deck central in the
mark-$D_4$; the clock's order itself is fermionically forced, $U^2=(-1)^F$ exactly
(\veri{v506\_seam\_clock\_rigidity.py}, \texttt{SEAM.CLOCK.RIGIDITY.01}; markers unchanged); and the
bit's origin is settled --- the deck's \emph{class} is derived (a half-period translation of the
seam double, equivalence tetrad with the nonsplit NS Fock face), only its \emph{position} stays the
carrier input (\veri{v507\_seam\_bit\_origin.py}, \texttt{SEAM.BIT.ORIGIN.01}; markers unchanged);
and that position half is now \emph{topology} --- the covering deck is free on the seam circle, the
edge class is excluded, and the residual is the square-modulus datum $\tau=i$ alone
(\veri{v510\_seam\_bit\_freedom.py}, \texttt{SEAM.BIT.FREEDOM.01}; markers unchanged) ---
equivalently, since \vref{v512}: one discrete symmetry-lift bit (flag transitivity of the four
marks, $V_4\to D_4$, $\iff\tau=i$). That bit has since survived \emph{ten} side-blind derivation
attacks --- free RP/$\Theta$ (the eighth, \vref{v521}), the exhausted Wick-computable
free-plus-twist state class (the ninth, \vref{v525}), the first genuinely interacting seam toy
(the tenth, \vref{v529}) --- and is now physically \emph{defined} as the twist-class choice:
a $15$-facet exact equivalence web with a gauge-robust order parameter ($O=\tfrac12$ at $m=4$,
$0$ elsewhere), in principle interferometrically readable (\vref{v528},
\texttt{SEAM.BIT.TWISTCLASS.01}; the bit remains formal input, no marker moves) --- and it now
carries the first \emph{positive} selection datum: reflection positivity, imposed on the
interacting mark family, keeps exactly \emph{one} member alive, $\delta=\pi/2$ with positive
coupling --- the carrier of that order parameter
(\veri{v534\_seam\_straddle\_cone.py}, \texttt{SEAM.STRADDLE.CONE.01}; toy-level evidence for a
dynamical origin, not a derivation).

\begin{proposition}[Per-sector assignment from the cusp ordering --- part (i) closed]
A parabolic line sub-bundle $\ell$ of the slope-$0$ semistable $E$ obeys
$\pardeg(\ell)=d_\ell+w^{\mathrm{cusp}}_\ell+(\text{other weights})\le 0$, so a \emph{larger} right-handed
cusp weight forces a \emph{more negative} degree $d_\ell$. Hence the sector with the largest cusp carries
the most-negative summand $\mathcal O(-2)$, i.e.\ the exponent $2$. With the $D_5$ cusps
\[
  \text{lepton}\leftrightarrow 2y_+=1\ >\ \text{up}\leftrightarrow -2y_-=\tfrac23\ >\ \text{down}\leftrightarrow 2(y_++y_-)=\tfrac13 ,
\]
the unique exponent $2$ goes to the lepton sector, giving
\begin{equation}
  \boxed{\ (a_u,a_d,a_e)=(1,1,2)=\bigl(\operatorname{round}(2\,\mathrm{cusp}_f)\bigr)_f\ }\qquad\mE/\mC ,
\end{equation}
the unique stable distribution of the multiset $\{1,1,2\}$ over the cusp-ordered sectors. Physically the
lepton's first generation (electron) is the lightest first-generation fermion, so it is the most
suppressed --- exponent $2$ --- consistent with $m_e<m_u<m_d$.
\end{proposition}

\noindent\textbf{Fused statement.} The anchors are now fully fixed: the \emph{multiset} $\{1,1,2\}$ from
$A_3$ stability (Schur--Horn $+$ balance) $=$ splitting type $\mathcal O(-2)\oplus\mathcal O(-1)^2$, and
the \emph{assignment} lepton$\to2$, up$\to1$, down$\to1$ from the $D_5$ cusp ordering via parabolic
stability. Multiset from $A_3$, assignment from $D_5$.
\par\smallskip\noindent\textbf{The single remaining piece --- part (ii).}
Only the \emph{full} residue sets $\{0,1,3\}$, $\{1,2,5\}$, $\{2,3,5\}$ stay open. The value $3$ exceeds
the exponent-at-infinity bound $\tfrac83$, so it is a genuine \emph{per-puncture} parabolic degree (an
asymmetric Hecke saturation at the individual punctures), not an exponent at infinity. The anchors (the
first-generation entries) are fixed above; the remaining two entries per sector require the non-symmetric
per-puncture saturation. \mC
\par\smallskip

\subsection{Part (ii): the $\mathbb Z_6$ structure of the residues, and the frontier}
Every residue decomposes \emph{uniquely} into a cusp index $c\in\{0,1,2\}$ (the $\mathbb Z_3$ family
phase, $c=3w$) and a sheet parity $s\in\{0,1\}$ (the $\mathbb Z_2$ deck involution), via the bijection
$r=2c+3s\bmod 6$:
\[
  0\!\to\!(0,0),\ 1\!\to\!(2,1),\ 2\!\to\!(1,0),\ 3\!\to\!(0,1),\ 4\!\to\!(2,0),\ 5\!\to\!(1,1).
\]
The full transport table in these parabolic coordinates is
\begin{center}\small
\begin{tabular}{l|ccc|ccc}
\toprule
 & \multicolumn{3}{c|}{cusp index $c$ (g1,g2,g3)} & \multicolumn{3}{c}{sheet $s$ (g1,g2,g3)} \\
\midrule
up     & $2,0,0$ & & & $1,1,0$ & & \\
down   & $2,1,1$ & & & $1,1,0$ & & \\
lepton & $1,1,0$ & & & $0,1,1$ & & \\
\bottomrule
\end{tabular}
\end{center}
\textbf{Provable structure \mE:} (a) windings are first-generation only, $n=(1,0,0)$ per sector,
$\sum n=\Nfam=3$; (b) every sector has exactly \emph{two} sheet-odd generations (sheet row-sums all
$=2$); (c) the gen-$1$ entries are the anchors $(1,1,2)$ of part~(i). The \emph{column} sums of the two
coordinate tables read carrier and Coxeter data:
\begin{equation}
\begin{aligned}
  &\text{cusp-index } C\!:\ \text{col sums }(5,2,1)=(\gcar,\,s,\,N_\Phi),\ \textstyle\sum=8=r(E_8);\\
  &\text{sheet } S\!:\ \text{col sums }(2,3,1)\ \text{(a perm.\ of }A_3\text{ exponents }1,2,3),\ \textstyle\sum=6=|R^+(A_3)| .
\end{aligned}\tag{$\ast$}
\end{equation}

\begin{keybox}[The $\mathbb Z_3\times\mathbb Z_2$ lift lemma --- the winding is almost automatic]
Assembling the cusp matrix $C$ and sheet matrix $S$ above into the natural integer lift
$T:=2C+3S$ gives
\[
  T=\begin{pmatrix}7&3&0\\7&5&2\\2&5&3\end{pmatrix},\qquad
  T\bmod 6=R,\qquad L-T=6\,E_{31},
\]
i.e.\ the $\mathbb Z_3\times\mathbb Z_2$ lift already produces the $+6$ overflow of the up- and
down-sector first generation; the \emph{only} extra global $+6$ winding sits on the lepton first
generation (the strongest cusp, $\mathcal O(-2)$). Moreover $\sum T=34=p_5(a)$ and $\sum L=40=p_1p_3$, so
the natural lift carries the anchor power $p_5$ and the single lepton winding completes it to the flavor
budget $40$. \mE
\end{keybox}

\subsection{The flavor transport theorem}
We now state the result as a theorem: under three precisely-stated hypotheses --- two of which are
already established in this note --- the entire transport matrix (anchors \emph{and} the six non-anchor
residues) is forced. The combinatorial core is proved by complete enumeration.

\begin{theorem}[Flavor transport residues on $\mathbb P^1\setminus\mu_4$]
\label{thm:flavor}
Let the three charged sectors $f\in\{u,d,e\}$ have transport residues $r_{f,j}\in C_6=\mathbb Z_6$
($j=1,2,3$). Assume:
\begin{itemize}[leftmargin=1.6em]
\item[\textbf{(H1)}] \emph{(non-degeneracy)} within each sector the three residues have pairwise
\emph{distinct} cyclic distances on $C_6$;
\item[\textbf{(H2)}] \emph{(cusp labelling)} the three sectors are the three $D_6$-images of one base
triple under the cusp action of $\mathbb P^1\setminus\mu_4$: $u=\mathrm{id}$, $d:x\mapsto 2{-}x$,
$e:x\mapsto 2{+}x$ (fixed by the right-handed cusps $-2y_-,\,2(y_+{+}y_-),\,2y_+$);
\item[\textbf{(H3)}] \emph{(anchor \& winding)} the first (lightest) generation of each sector carries
the single full $+6$ winding, with anchor residue $(a_u,a_d,a_e)=(1,1,2)$, and within a sector lengths
decrease with generation (mass hierarchy).
\end{itemize}
Then the base triple is uniquely $\{0,1,3\}$, the sector residue sets are
$\{0,1,3\},\{1,2,5\},\{2,3,5\}$, and the full transport matrix is
\begin{equation}
  \boxed{\ L_u=(7,3,0),\quad L_d=(7,5,2),\quad L_e=(8,5,3)\ }\qquad
  \Bigl(\textstyle\sum L=40,\ (G_1,G_2,G_3)=(22,13,5)\Bigr).
\end{equation}
\end{theorem}

\begin{proof}
\emph{(Step 1 --- base triple.)} By complete enumeration of the $\binom63=20$ three-subsets of
$\mathbb Z_6$, the cyclic-distance multiset is one of $\{1,1,2\}$ (six subsets), $\{2,2,2\}$ (two), or
$\{1,2,3\}$ (twelve). Hypothesis (H1) selects the distinct case $\{1,2,3\}$. These twelve subsets form a
single $D_6$-orbit with trivial stabiliser (verified: the only $(r,s)$ with $sx+r\equiv\{0,1,3\}$ is the
identity), of which $\{0,1,3\}$ is the representative. \emph{(Step 2 --- sets.)} Applying the three
$D_6$-elements of (H2) gives $u:\{0,1,3\}$, $d:2{-}\{0,1,3\}=\{1,2,5\}$, $e:2{+}\{0,1,3\}=\{2,3,5\}$;
the trivial stabiliser guarantees these are three \emph{distinct} sets, one per sector.
\emph{(Step 3 --- matrix.)} By (H3) the anchor occupies generation $1$ with a $+6$ winding; deleting it
from the set leaves exactly two residues, which the monotone order assigns to generations $2,3$:
\[
  u:\{0,1,3\}\!\setminus\!\{1\}=\{0,3\}\Rightarrow(1{+}6,3,0),\quad
  d:\{1,2,5\}\!\setminus\!\{1\}\Rightarrow(7,5,2),\quad
  e:\{2,3,5\}\!\setminus\!\{2\}\Rightarrow(8,5,3).
\]
Summing gives $\sum L=40$ and column sums $(22,13,5)$.
\end{proof}

\noindent\textbf{Status of the hypotheses (honest grounding).}
\begin{itemize}[leftmargin=1.6em]
\item (H1) is the parabolic genericity/spread statement (distinct masses $\Rightarrow$ distinct
transport gaps); it is physically forced but stated here as a hypothesis. \mC
\item (H3) is \emph{derived}: the anchor multiset $\{1,1,2\}$ is the Schur--Horn$+$stability splitting
type (Section above), and the per-sector assignment $(1,1,2)$ is the $D_5$ cusp ordering. \mE/\mC
\item (H2) is now \emph{small and not a free choice}: with (H1),(H3) the up and lepton sectors are fixed
(sum $4$ unique; sum $10$ with anchor $2$ unique), leaving only the down branch $\{1,2,5\}$ vs
$\{1,3,4\}$. The transport-kernel theorem \emph{derives} the former (down $=\{1,2,5\}$) from the
uniquely-fixed family-connection holonomy, and it satisfies the spectral selector ($\det R=h(D_5)$,
principal minors $(2,3,5)$); the sibling $\{1,3,4\}$ does not arise. So (H2)'s answer is determined; only
the parabolic$\leftrightarrow$transport equivalence remains. \mE/\mC
\end{itemize}
So the matrix is a \emph{theorem} (combinatorial core fully proved), resting on three hypotheses of which
(H3) is derived, (H1) is a genericity statement, and (H2) is the single cusp-geometry input.

% =====================================================================
